% ============================================================================
%  Subdocumento 14. Contenido. Se usa igual en el documento único y en el
%  PDF propio del subdocumento. Los formularios que le asigna el T-21 se
%  agregan solos al final (datos en formularios/ de esta carpeta).
% ============================================================================

\begin{resumenApertura}
\marcador{Qué resuelve este subdocumento, en cuatro o cinco líneas}
\begin{recibe}
  \item \marcador{Qué recibe el mandante}
\end{recibe}
\end{resumenApertura}

\section{Propuesta de valor}\label{sec:14-propuesta-de-valor}

\marcador{Síntesis desde la ingeniería integral}

\section{Beneficios cuantitativos para el mandante}\label{sec:14-beneficios-cuantitativos-para-el}

\marcador{Desempeño, restauración, disponibilidad y ahorro operacional, sin montos}

\section{Equilibrio entre alcance, tiempo, costo y calidad}\label{sec:14-equilibrio-entre-alcance-tiempo-}

\marcador{Cómo se equilibran}

\section{Coherencia arquitectónica y tecnológica}\label{sec:14-coherencia-arquitectonica-y-tecn}

\marcador{Entre todas las secciones}

\section{Trazabilidad de extremo a extremo}\label{sec:14-trazabilidad-de-extremo-a-extrem}

\marcador{De requerimiento a operación}
